\section{Discussion}
This work demonstrates that evolutionary approaches for designing agentic systems as a viable and effective paradigm to AI safety research. \method~addresses fundamental limitations of manually designed red-teaming systems by enabling the rapid discovery of effective system designs tailored to specific target models and evaluation benchmarks. Across multiple benchmarks and target models, \method~produces systems that achieve new state-of-the-art performance. At the same time, \method~highlights the importance of human expertise in scientific discovery, serving both as a strong initial baseline and as guidance for automated exploration. As companies and organizations increasingly train, deploy, and integrate AI systems into critical domains, the demand for robust, adaptive safety evaluation grows accordingly. \method~offers a scalable solution for designing highly effective red-teaming systems that can uncover potential safety vulnerabilities without constant manual intervention. We acknowledge several limitations of \method~along with promising directions of future work below.

\label{appx: limitation & future works}
%the resulting adversarial prompts are either not human-readable (Zou et al., 2023), hence they can be easily filtered by perplexity-based mitigation strategies (Jain et al., 2023), or require computationally intensive discrete optimization over the combinatorial token space to generate a single adversarial promp
\textbf{Self-evolving Algorithms.} Evolutionary algorithms, though a concept with over two centuries of history, have regained prominence in the recent automated scientific discovery landscape, with applications in high-value domains such as algorithm optimization, GPU kernel engineering, and LLM training \citep{novikov2025alphaevolvecodingagentscientific,yuksekgonul2026learningdiscovertesttime,karpathy2026autoresearch}. These methods all require a well-contained task with a mature interface and verifiable reward signal, which is precisely the framework that \method~provides. We hope \method~serves as a pioneering step toward broader adoption of more sophisticated self-evolving algorithms in AI safety research.

\textbf{Attacker-Target Co-Evolution.} In practice, target models are not static, as safety teams continuously patch discovered vulnerabilities. This creates a co-evolving arms race between red-teaming systems and defensive measures. Future research could explore co-evolutionary frameworks where both attacking and defending models evolve simultaneously, potentially yielding more robust security insights.

\textbf{Addressing Systematic Mode Collapse.} We observe evidence of mode collapse \citep{jiang2025artificialhivemindopenendedhomogeneity} in the design space exploration. Because meta agents are trained on literature corpora containing existing red-teaming research, they tend to converge on similar design approaches across generations. For example, the discovered systems in Mixtral attacker targeting Llama2-7B pipeline consistently incorporate recurring techniques such as attack templates (used 10 times over 10 generations), refusal suppression (8 times over 10 generations), and genetic operators (mutation and crossover, 10 times over 10 generations). We further quantify the similarity across different generations using their semantic embedding in Appendix \ref{appx: similarity across generated systems}. While these designs represent strategies validated by prior research, their prevalence limits exploration of the broader design space, preventing the discovery of novel attacks. To encourage diversity and originality, we need to address the mode collapse phenomenon that arises during the evolutionary stage.

%\textbf{Improving Reproducibility}. Our experiments rely on several closed-source models as meta agent, target model and judge model. The ablation study using a different meta agent LLM (Appendix \ref{appx: deepseek meta agent}) confirms that performance depends on the meta agent's reasoning capability. Since \method uses an inaccessible proprietary model (gpt-5-2025-08-07) as the meta agent, reproducibility is limited. Future work can explore open-weight models as meta agents to improve reproducibility.

%\textbf{Improving Query Time.} \method~incurs substantial computational overhead during the search. Generating multiple candidate systems per generation, each requiring initial evaluation, demands significant inference-time compute. As shown in Appendix \ref{appx: query per success}, while the resulting systems achieve strong performance with reasonable test-time efficiency (339 queries per success), the training-time cost is considerably higher. In particular, when accounting for all queries used to evaluate candidate systems across generations, \method~requires a large number of training-time queries (122k), which includes the queries needed to evaluate all the systems produced for each generation until the best performing system is produced. Future works can explicitly incorporate query efficiency as an additional objective for optimization.

%\textbf{Pushing the Quality-Diversity Frontier.} Our diversity analysis uncovered a critical trade-off: optimization pressure for ASR can drive convergence toward homogeneous attack patterns. The reward shaping and $\epsilon$-constraint sample rejection method shows promise in pushing the quality-diversity frontier, but introduces additional hyperparameters (i.e. rejection threshold, or $\epsilon$). More sophisticated multi-objective optimization techniques may be applied for further study. In science practice, researchers must often balance multiple competing objectives, such as cost, latency, originality and performance. Therefore, further research on multi-objective optimization is promising.




%\textbf{Originality} As evidenced in the limitations above, we have observed homogeneous design pattern across different target models and benchmarks. To achieve diversity and originality in scientific discovery, we need to look into ways to address the mode collapse phenomenon through methods such as post-training conditioning.

%\textbf{Explicit reinforcement learning formulation} \method~implicitly incorporates reinforcement learning through its structural guidance that learns from higher-performing systems. Reformulating this as an explicit reinforcement learning problem could accelerate convergence and improve sample efficiency.